%! Modeling with Quadratics — Unit 2, Lesson 2.7 — Student Packet Cover Sheet
\documentclass[10pt]{article}
\usepackage{algebra2-article}
\usepackage{algebra2-boxes}
\usepackage{ltablex}
\keepXColumns

\begin{document}

% Full-bleed forest banner
\begin{tikzpicture}[remember picture, overlay]
  \fill[forest] (current page.north west)
    rectangle ([yshift=-0.9in]current page.north east);
\end{tikzpicture}
\vspace{-0.8in}

\begin{center}
  {\color{white}\textbf{\LARGE Algebra 2}}\\[4pt]
  {\color{forestmid}\large\textbf{Unit 2 \quad Quadratic Functions}}\\[6pt]
  {\color{white}\textbf{\large Lesson 2.7 \quad Modeling with Quadratics}}
\end{center}

\vspace{0.22in}
\namedateperiod
\vspace{0.18in}

\begin{learningtargetbox}
\small
\begin{enumerate}[leftmargin=1.5em, itemsep=5pt]
  \item \ldots{} \textbf{build} a quadratic model from a \textbf{projectile}, \textbf{maximum-area}, or \textbf{revenue} situation, defining my variable with units.
  \item \ldots{} read the answer off the model's key features --- the \textbf{$y$-intercept} as the starting value, a \textbf{zero} as ``reaches $0$,'' and the \textbf{vertex} as the \textbf{maximum or minimum}.
  \item \ldots{} \textbf{interpret} every answer in context with units and \textbf{reject} values that don't fit the situation (negative time, an impossible width).
\end{enumerate}
\end{learningtargetbox}

\vspace{0.14in}

\begin{tocbox}
\small
\renewcommand{\arraystretch}{1.5}
\begin{tabularx}{\linewidth}{@{} c l X r @{}}
\rowcolor{forestbg}
\textbf{\#} & \textbf{Component} & \textbf{Description} & \textbf{Score} \\
\midrule
1 & Warm-Up        & Spiral review: find a vertex, solve a quadratic, read a height model & \blank{1.2cm} \\
2 & Guided Notes   & Projectile, maximum area, and revenue models --- and the feature-to-question map & \blank{1.2cm} \\
3 & Group Activity & \emph{Model It!} --- read, build, and interpret quadratic models, by tier        & \blank{1.2cm} \\
4 & Exit Ticket    & Quick check: max height \& landing, a max-area model, plus an SOL item          & \blank{1.2cm} \\
5 & Homework       & Build \& interpret projectile / area / revenue models; find the error           & \blank{1.2cm} \\
\midrule
  & \multicolumn{2}{r}{\textbf{Total}} & \blank{1.2cm} \\
\end{tabularx}
\end{tocbox}

\vspace{0.10in}

\begin{remindbox}
\small To \textbf{model} with a quadratic: \textbf{(1)} define your variable (with units), \textbf{(2)} write the
quadratic, \textbf{(3)} decide which \emph{feature} the question wants, \textbf{(4)} find it. The \textbf{$y$-intercept}
$c$ is the \emph{starting value}; a \textbf{zero} (positive root) is \emph{when/where it reaches $0$}; the
\textbf{vertex} is the \textbf{maximum or minimum} --- its input tells you \emph{when/where}, its output tells you
\emph{how much}. Find the vertex with $x=-\frac{b}{2a}$, then evaluate. Always \textbf{interpret in context with units}
and throw out answers that can't happen (negative time, a width outside the fence). The projectile model is
$h(t)=-16t^2+v_0t+h_0$.
\end{remindbox}

\end{document}
